\section{Integration decisions and resulting usefulness}
\label{sec:want-completion}
This subversion integrates definitions where their types and physical domains
can be stated explicitly. It keeps earlier corrected claims corrected and
separates source compatibility from a changed exact algebraic profile.

\small
\begin{longtable}{@{}P{13mm}P{43mm}P{94mm}@{}}
\toprule ID & Source contribution & R12 decision and reason\\\midrule\endhead
R12-01 & W1 p18, W0 node word & Preserve all W64 fields separately from JK32 and the 512-bit WI record; equal bit width does not establish equal meaning.\\
R12-02 & W0 pinion and OTAN2 & Name discrete rounded and exact algebraic pinions separately; keep finite phase permutation separate from UGTS ratio/arctangent.\\
R12-03 & W1 p18, determinant & Derive the singular values and scaled passive dilation; volume preservation alone does not imply passive optical power preservation.\\
R12-04 & W1 p5, event feedback & Completed observations affect the next command; retain an explicit delay and consume each record once in the optional JK event lane.\\
R12-05 & W1 p19, transport & Preserve raw bytes and bind calibration, clock epoch, drift and uncertainty outside the packet; no inferred GPS chronology from a bare counter.\\
R12-06 & W1 pp6--7, rim & Retain local section and conditional compliance results; supply a finite modal weak form rather than treating a PDE as an implemented finite ODE.\\
R12-07 & W1 pp8--10, gauge & Retain exact bridge sign/inverse, arm powers and physical root domain; carry strain transfer and temperature calibration.\\
R12-08 & W1 pp11--12, optics & Retain bounded scalar diffraction and passive energy constraints; the slit is an aperture unless a waveguide is separately specified.\\
R12-09 & W1 p12, two bands & Counts follow calibrated spectral routes; one photon is not required to appear in both nonoverlapping bands.\\
R12-10 & W1 pp13,16, memory & Preserve reporter and thermal expressions before rounding; add zero-rate/conductance branches and held-input hypotheses.\\
R12-11 & W1 p14, gate & Condition false-acceptance bounds on the same validity/timing population; two channels do not imply independence.\\
R12-12 & W1 p15, instrument & Represent finite complex operators through real pairs and retain outcome probabilities; no finite registry claim for an unrestricted infinite-mode system.\\
R12-13 & W1 p17, PSI & Extend exclusion to moving slit geometry with explicit bounds; optical diffraction and heat need separate retained terms or error budgets.\\
R12-14 & W1 p20, calibration & Preserve monotone inverse intervals and identifiability requirements; exact nominal coefficients do not imply exact measurements.\\
R12-15 & W1 p23, calculations & Retain historical reports as source evidence only; no new algorithm test, calculation example or physical experiment was run.\\
R12-16 & Photonic execution & Add a conditional refinement, margin and finite-workload reliability contract; no complete backend follows from a transfer equation alone.\\
\bottomrule
\end{longtable}
\normalsize

\subsection{Useful physical, digital and spatial applications}
The combined equations now describe a deforming optical aperture with strain
readout, a calibrated optical-event receiver, thermal/material memory, and a
causal controller connected to the existing word kernel. The same descriptions
support instrument geometry, alignment monitoring, pulse timing, bounded
geometric exclusion and lossless digital observation replay. An exact expression
record can preserve how an event was derived from its inputs and calibration.
It cannot retroactively measure an absent channel.

For a ground-station optical head, the geometry and strain operators can model
mount displacement or aperture deformation; the optical equations can model
collection and detection; and the clock adapter can associate observations with
the station's declared time scale. Those are useful instrument-level connections
to the retained satellite/station coordinate model. They do not add an orbital
force law, improve an ephemeris, remove atmospheric or timing error, or establish
a new MEO/GPS prediction tolerance. R10's dated numerical results keep their
original data, horizon and implementation scope.

For photonic computation, the scaled pinion is a specific ideal linear-optical
candidate and the bit-margin contract gives a precise route to a digital
correctness claim. These are stronger design statements than treating a
geometric drawing or a conserved determinant as a physical proof. A complete
device still needs the encoding, memory, nonlinear decisions, restoration,
timing, resources and component response required by that contract.

\subsection{Remaining inputs are specific, not an invented tolerance}
\begin{center}\small
\begin{tabularx}{\linewidth}{@{}P{35mm}Y@{}}
\toprule Intended claim & What is still needed to establish it\\\midrule
Mechanical displacement & Chosen specimen, supports, dimensions, material tensor, loads and adequate basis/model scope\\
Optical collection & Wavelengths, refractive indices, aperture/lens geometry, measured or justified transfer and collection loss\\
Physical event meaning & Detector response, reporter coupling, calibration, timing and joint noise statistics for the chosen mode\\
Thermal behavior & Capacitances, conductances, electrical/optical power and energy partition over the operating interval\\
Photonic exact executor & Explicit device architecture and one-step refinement over a declared finite resource and disturbance envelope\\
System tolerance & A specified output quantity and admitted input/model/error envelope, propagated through the relevant equations\\
\bottomrule
\end{tabularx}\end{center}
The ``functional tolerances of the system'' are consequently conditional bounds
such as the bridge domain, gate margin, timing window and geometric exclusion
inequality. They become numerical device limits only after their parameters
are supplied and justified. Lossless packing does not assign those parameters.

\subsection{Review method and preserved lineage}
The new mathematics was reviewed by derivation and independent source reading.
The proof register covers namespace/word preservation, causal lane composition,
passive pinion dilation, aperture contraction, finite instrument closure,
modal positivity, compliance, bridge inversion, reporter bounds, event
conditioning, moving geometry, passive heat and hardware refinement.
These are conditional arguments, not proof-assistant certificates or measured
performance. The rendered PDF is inspected for legibility separately.

The following R11 chapters and complete inherited R10 corpus preserve their
source bytes. Their historical release labels and performance statements are
retained inside clearly marked inherited sections. R12 does not rerun the
source's programs or reclassify their outputs as new evidence. The release
includes the editable source, provenance, review records and distribution
manifest so each adopted statement can be traced to its source or derivation.

\begin{panel}
\textbf{Result.} WANTWOMBAN materially helps formalize an optical sensing and
feedback interface and supplies an ideal passive candidate for a scaled pinion.
The exact seed can retain the relevant operators and input lineage losslessly.
A functioning photonic exact processor and its physical accuracy remain
implementation and characterization claims that this mathematical integration
does not yet establish.
\end{panel}
